\documentclass[11pt,a4paper]{article}
\usepackage{gtu-el-yazisi}
\begin{document}
\notbilgi{Alan Teorisi}{Final çalışması: Çıkmış sınav 2023}{Cikmis-2023.pdf}
\textit{Not: Orijinal dosyadaki basılı sayfalar (s. 1) dijitalleştirilmedi; yalnızca el yazısı sayfalar aktarıldı.}

\sayfa{2}
\section{Soru 1}
\sekilnotu[2]{Sayfanın üstünde basılı sınav sorusu (Soru 1, şıklar a, b, c) var; dijitalleştirilmedi.}

\elcizimi[0.5]{sekiller/final-calismasi-Cikmis-2023-s2-1.png}[2]

\elcizimi[0.4]{sekiller/final-calismasi-Cikmis-2023-s2-2.png}[2]

\[ \rho_v(x,y,z)=\rho_c\,\delta(x)\,\delta(y)\,\bigl[u(z+\ell/2)-u(z-\ell/2)\bigr] \]
\[ Q=\ell\cdot\rho_c \qquad\qquad \rho_c=\frac{Q}{\ell} \]
\[ \Rightarrow\ \rho_v(x,y,z)=\frac{Q}{\ell}\,\delta(x)\,\delta(y)\,\bigl[u(z+\ell/2)-u(z-\ell/2)\bigr] \]

b)
\[ \oint_S \vec D\cdot d\vec s = Q \]
\[ \oint_S \varepsilon_0\,\vec E\cdot d\vec s = Q=\ell\rho_c \]
\[ \int_{z=-\ell/2}^{\ell/2}\int_{\phi=0}^{2\pi}\varepsilon_0\,E_\rho(\rho)\,\underbrace{\rho\,d\phi\,dz}_{ds}=\ell\rho_c=Q \]
\[ 2\pi\ell\,\varepsilon_0\,\rho\,E_\rho(\rho)=\ell\rho_c \]
\[ E_\rho(\rho)=\frac{\rho_c}{2\pi\varepsilon_0\rho}\ \Rightarrow\ \vec E=\frac{\rho_c}{2\pi\varepsilon_0\rho}\,\vec e_\rho \]

\sayfa{3}
ve $\rho=\sqrt{x'^2+y'^2}$ şeklindedir
\[ E(0,y,0)=\frac{\rho_c}{2\pi\varepsilon_0\sqrt{\okunamadi{}\,y^2}}=\frac{\rho_c}{2\pi\varepsilon_0\,y}=\frac{Q}{2\pi\varepsilon_0\,y\,\ell}\qquad\Bigl(\rho_c=\frac{Q}{\ell}\Bigr) \]
\[ E(0,d,0)=\frac{Q}{2\pi\varepsilon_0\,d\,\ell} \]

c) P noktasında konulan $q$ yüküne etkiyen kuvvet
\[ \vec F_q=q\cdot E(0,d,0)=q\cdot\frac{Q}{2\pi\varepsilon_0\,d\,\ell}\ \text{(N) kadar olup $L$'nin} \]
\noindent $q$'ya uyguladığı kuvvettir.

\elcizimi[0.5]{sekiller/final-calismasi-Cikmis-2023-s3-1.png}[3]

\begin{align*}
d\vec F &= \frac{q}{4\pi\varepsilon_0}\cdot\frac{dQ}{r^2}\,\vec e_r\\
&= \frac{q}{4\pi\varepsilon_0}\,\frac{\rho_v\,dv}{r^2}\,\vec e_r
\end{align*}
\[ \vec F=\frac{q}{4\pi\varepsilon_0}\,\frac{Q}{\ell}\int_{r=0}^{\okunamadi{}}\int_{\phi=0}^{2\pi}\int_{\theta=0}^{\pi}\frac{1}{r^2\footnotemark}\,r^2\,\frac{\vec r}{r}\,\sin\theta\,dr\,d\phi\,d\theta \]
\footnotetext{Orijinalde paydaki $r^2$ çizilerek iptal edilmiş.}
\[ d\vec F=\int_0^\pi\frac{q\cdot\dfrac{Q}{\ell}}{4\pi\varepsilon_0\,d}\cdot(\sin\theta) \]
\[ \cos\pi=-1 \qquad \cos 0=-1\,\footnotemark \]
\footnotetext{Claude notu: muhtemelen $\cos 0=1$ olmalı.}
\[ \vec F_Q=-\frac{q\cdot Q/\ell}{2\pi\varepsilon_0\,d\,\ell} \]

\sayfa{4}
\section{Soru 2}
\sekilnotu[4]{Sayfanın üstünde basılı sınav sorusu (Soru 2, şıklar a, b) var; dijitalleştirilmedi.}

\elcizimi[0.7]{sekiller/final-calismasi-Cikmis-2023-s4-1.png}[4]

\[ z>0\qquad \Delta V=\frac{-\rho_v}{\varepsilon_1}=\frac{-q}{\varepsilon_1}\,\delta(x-\overset{0}{x'})\,\delta(y-\overset{0}{y'})\,\delta(z-\overset{h}{z'}) \]
\[ =\frac{-q}{\varepsilon_1}\,\delta(x)\,\delta(y)\,\delta(z-h) \]

\sayfa{5}
\section{Soru 3}
\sekilnotu[5]{Sayfanın üstünde basılı sınav sorusu (Soru 3, şıklar a, b) ve Şekil-1 ($Oyz$ düzleminde $z$ yönünde $I_1$, $I_2$, $I_3$ akımı taşıyan üç paralel tel; aralarındaki uzaklıklar $a$ ve $b$) var; dijitalleştirilmedi.}

\[ \vec B=\frac{\mu_0}{2\pi\rho}\,I\,\vec e_\phi \qquad\text{(Biot-Savart)} \]
$\vec B$ altında bulunup üzerinde $I$ akımı akan bir tele etki eden kuvvet
\[ d\vec F=I\,d\vec\ell\times\vec B \]

\[ d\vec F_2=I_2\,dz\,\vec e_z\times\vec B_2 \]
\[ \vec B_2=\frac{\mu_0}{2\pi}\Bigl[\frac{I_3}{b}-\frac{I_1}{a}\Bigr]\vec e_x \]
\begin{align*}
\vec F_2 &= \int_0^{\ell}d\vec F_2=\int_0^{\ell}I_2\,\vec e_z\times\vec B_2\,dz\\
&= \int_0^{\ell}\frac{I_2\mu_0}{2\pi}\Bigl[\frac{I_3}{b}-\frac{I_1}{a}\Bigr]\underbrace{\vec e_z\times\vec e_x}_{\vec e_y}\,dz
\end{align*}
\[ \vec F_2=\frac{I_2\mu_0}{2\pi}\Bigl[\frac{I_3}{b}-\frac{I_1}{a}\Bigr]\vec e_y\ \text{olarak bulunur} \]

\par\noindent\rule{\linewidth}{0.4pt}\par

\[ \vec B_1=\frac{\mu_0}{2\pi}\Bigl[\frac{I_3}{a+b}+\frac{I_2}{a}\Bigr]\vec e_x \]
\[ d\vec F_1=I_1\,dz\,\vec e_z\times\vec B_1 \]
\begin{align*}
\vec F_1 &= \int_0^{\ell}d\vec F_1=\int_0^{\ell}I_1\,\vec e_z\times\vec B_1\,dz\\
&= \int_0^{\ell}\frac{I_1\mu_0}{2\pi}\Bigl[\frac{I_3}{a+b}+\frac{I_2}{a}\Bigr]\underbrace{\vec e_z\times\vec e_x}_{\vec e_y}\,dz
\end{align*}

\sayfa{6}
\[ \vec F_1=\frac{I_1\mu_0}{2\pi}\Bigl[\frac{I_3}{a+b}+\frac{I_2}{a}\Bigr]\vec e_y \]

\par\noindent\rule{\linewidth}{0.4pt}\par

\[ \vec B_3=\frac{\mu_0}{2\pi}\Bigl[-\frac{I_1}{a+b}-\frac{I_2}{b}\Bigr]\vec e_x \]
\elcizimi[0.18]{sekiller/final-calismasi-Cikmis-2023-s6-1.png}[6]

\[ d\vec F_3=I_2\,dz\,\vec e_z\times\vec B_3\,\footnotemark \]
\footnotetext{Claude notu: muhtemelen $I_3$ olmalı (sonraki satırlarda da $I_2$ yazılmış).}
\begin{align*}
\vec F_3 &= \int_0^{\ell}-I_2\,\frac{\mu_0}{2\pi}\Bigl[\frac{I_1}{a+b}+\frac{I_2}{b}\Bigr]\underbrace{\vec e_z\times\vec e_x}_{\vec e_y}\,dz\\
&= -\frac{I_2\mu_0}{2\pi}\Bigl[\frac{I_1}{a+b}+\frac{I_2}{b}\Bigr]\vec e_y
\end{align*}

\section{Silindirik}
\[ x=\rho\cos\phi \qquad y=\rho\sin\phi \]
\[ \vec r=\rho\cos\phi\,\vec e_x+\rho\sin\phi\,\vec e_y+z\,\vec e_z \]
\[ \int_{z=0}^{\okunamadi{D}}\int_{\phi=0}^{2\pi}\int_{\rho=0}^{a}dV=\rho\,d\rho\,d\phi\,dz \]

\[ d\vec F=\rho_s\,ds\,\vec E \]
\[ ds=a^2\sin\theta\,d\phi\,d\theta \]
\[ \vec F=\int_{\theta=0}^{\pi}\int_{\phi=0}^{2\pi} \]

\section{Sınır koşulları}
\[ \rho_s=\varepsilon_1E_n^{(2)}-\varepsilon E_n^{(1)} \]
\[ \oint_S\vec D\cdot d\vec s=Q \]
\[ \vec E=\frac{Q}{4\pi\varepsilon_0 r^2}\,\vec e_r\ \ \text{genel} \]
Sonsuz uzun çizgisel yük dağılımı için
\[ \vec E=\frac{\rho_c}{2\pi\varepsilon_0\rho}\,\vec e_\rho \]

iletkenlerin

\end{document}
